\begin{abstract}
Regional English accent is a source of uneven automatic speech
recognition performance, but its effect on final speech-translation quality is
less clearly established. This survey examines how accent may affect Direct
speech translation, which maps speech to translated text without an exposed
transcript, and Cascaded speech translation, which connects speech recognition
to machine translation. Twenty-seven works were retained through thematic
selection covering system architectures, accent robustness, error
propagation, datasets, translation metrics, and statistical comparison. The
review finds that neither architecture is consistently superior: reported
outcomes vary with language pair, training resources, test data, references,
and metrics. Accent-related recognition disparities are well established, and
recognition errors can damage downstream translation, but these findings do
not alone demonstrate accent-related translation degradation. The closest
speech-translation studies either lack controlled regional-accent labels, omit
a matched Cascaded comparator, or evaluate an adjacent downstream task.
Common Voice and CoVoST 2 enable a relevant evaluation, although self-reported
accent labels, repeated content, speaker dependence, and reference design
constrain interpretation. The gap is a matched evaluation of Direct
and Cascaded systems on the same regional-English utterances using final
translation quality, motivating the question: how does regional English accent
affect the translation quality of these two system categories?
\end{abstract}

\begin{IEEEkeywords}
speech translation, regional English accents, robustness, direct speech
translation, cascaded speech translation, automatic speech recognition
\end{IEEEkeywords}
